\documentclass[10pt,a4paper]{amsart}
\usepackage[margin=19mm]{geometry}
\usepackage{amsmath,amssymb,mathtools}
\usepackage[scr=rsfs,cal=euler]{mathalfa}
\usepackage{xfrac}
\usepackage[unicode,hidelinks,bookmarks=false]{hyperref}
\newcommand{\ie}{i.e.\ }
\newcommand{\cf}{cf.\ }
\newcommand{\resp}{respectively\ }
\newcommand{\Subsection}{\subsection}
\providecommand{\nobreakdash}{\nobreak\hbox{-}}
\setlength{\emergencystretch}{2em}
\multlinegap0pt
\usepackage{bm}
\usepackage{bbm}
\usepackage{dsfont}
\usepackage{xspace}
\usepackage{cancel}
\usepackage{mathtools}
\usepackage[shortlabels]{enumitem}
\usepackage{amsthm}
\makeatletter
\@ifundefined{dummy}{\newtheorem{dummy}{Dummy}}{}
\@ifundefined{proposition}{\newtheorem{proposition}[dummy]{Proposition}}{}
\makeatother
\DeclareMathOperator{\id}{id}
\DeclareMathOperator{\image}{image}
\newcommand{\calC}{\mathcal{C}}
\newcommand{\calP}{\mathcal{P}}
\newcommand{\calT}{\mathcal{T}}
\title[Filtered complexes and cohomologically equivalent subcomplex]{Filtered complexes and cohomologically equivalent subcomplexes}
\author{Erlend Grong, Francesca Tripaldi}
\date{Source: arXiv:2308.11353v2 (CC BY 4.0)}
\begin{document}
\maketitle

\noindent\emph{Source and licence.} Filtered complexes and cohomologically equivalent subcomplexes. Version arXiv:2308.11353v2 (CC BY 4.0), distributed under the Creative Commons Attribution 4.0 International licence, as recorded in the arXiv licence metadata for this exact version and shown on the abstract page https://arxiv.org/abs/2308.11353v2. Full rights notice: licenses/arxiv-2308.11353-CC-BY-4.0.txt.\par\smallskip

\noindent\textit{Editorial scope.} Counted unit: the Proposition whose source label is \texttt{prop:Invariance} in the article (source0.tex lines 373--378, source label \texttt{prop:Invariance}), delivered together with the article's own complete original proof of it (source0.tex lines 379--406, one \texttt{proof} environment of 28 lines). No mathematics is written, repaired or re-derived: statement and proof are the article's own text line for line. Required context retained uncounted: Poperator (lines 231--231); Pi0operator (lines 327--327); prop:TransInv (lines 352--354). Every definition, display and label of those lines is retained unchanged. Omitted from this package and not redistributed: 1--230, 232--326, 328--351, 355--372, 407--1135. Nothing inside a retained excerpt is omitted or altered: every retained body is delivered exactly as the article's lines are written, so no omission receipt of any kind accompanies this package. Citations: the retained excerpts carry no citation command at all, so no bibliography entry of the article is transcribed and no reference is redistributed. Reference adaptation: none was needed and none was made; every reference inside the delivered unit resolves inside it, so no unresolved reference or citation is produced. Printed slips kept literal: no printed word, symbol, hypothesis or number of the counted statement or of its proof was altered. Layout and class: the article's own class is \texttt{amsart} with packages including lineno, fullpage, amsfonts, amsmath, amssymb, amsthm, pict2e, hyperref, amsbsy, xcolor, graphicx, enumerate, multicol, mathrsfs; the wrapper is the lane's pinned \texttt{amsart} editorial wrapper at 10pt on A4, which restates the theorem-like environments the retained text needs and adds the article's own macro definitions that the retained text needs (\texttt{\textbackslash calC}, \texttt{\textbackslash calP}, \texttt{\textbackslash calT}, \texttt{\textbackslash id}, \texttt{\textbackslash image}), with extra wrapper packages (bm, bbm, dsfont, xspace, cancel, mathtools, enumitem). The delivered numbering is the unit's own. Assets: no retained span contains a figure environment, a graphics-inclusion command, a PDF-page inclusion, a table or a TikZ picture; no image file is copied into this package, the package declares no asset dependency and no image file is redistributed with it. Licence: version arXiv:2308.11353v2 of this article is distributed under the Creative Commons Attribution 4.0 International licence, as recorded in the arXiv licence metadata for this exact version and shown on the abstract page https://arxiv.org/abs/2308.11353v2. A full rights notice is delivered at licenses/arxiv-2308.11353-CC-BY-4.0.txt. Characters: every retained excerpt is pure ASCII, so the delivered engine, XeLaTeX with its default Latin Modern text font, reports no missing character.
\begin{equation} \label{Poperator} P = \id - d_0^{-1} d -  dd_0^{-1}.\end{equation}

\begin{equation} \label{Pi0operator} \Pi_{0} = \id - d_0d_0^{-1} - d_0^{-1}d_0,\end{equation}

\begin{proposition} \label{prop:TransInv}
If $P^\infty = P^{S}$ is defined as in \eqref{Poperator} with respect to the base transversal and cotransversal $\calT \subset \calC$, while $\tilde \Pi_0$ is defined as in \eqref{Pi0operator} with respect to $\tilde \calT \subset \calC$ then $P^\infty\circ\tilde\Pi_0|_{\tilde \calP}$ is an isomorphism between $\tilde \calP$ and $\calP$ as filtered differential complexes.
\end{proposition}

\begin{proposition} \label{prop:Invariance}
\begin{enumerate}[\rm (a)]
    \item $\calT \subset \calC$ are also base transversals and cotransversals for $\tilde d_0$.
    \item If $P^\infty = P^{S}$ is defined as in \eqref{Poperator} with respect to $d_0$ and $\calT \subset \calC$ and $\tilde \Pi_0$ is defined as in \eqref{Pi0operator} with respect to $\tilde d_0$ and $\tilde \calT \subset \tilde \calC$, then $P^\infty\circ\tilde\Pi_0|_{\tilde \calP}$ is an isomorphism of differential complexes between $\tilde \calP$ and $\calP$.
\end{enumerate}
\end{proposition}

\begin{proof}
\begin{enumerate}[\rm (a)]
    \item Let us introduce the notation $\tilde d_0 = d_0 - N$, where the map $N = d_0- \tilde d_0 =  ( d-\tilde d_0) - (d- d_0)$ maps $K_j$ into $K_{j+1}$. It follows that $\hat N := d_0^{-1} N$ is nilpotent with $\hat N^{S+1} =0$, meaning that $A : =\id - \hat N$ is invertible.  If we consider the projection map $d_0d_0^{-1}$ applied to the image of $\tilde d_0$, we have that for any $\alpha\in K$
    \begin{align*}
        d_0d_0^{-1}(\tilde d_0\alpha)=d_0d_0^{-1}(d_0\alpha-N\alpha)=d_0\alpha-d_0\hat N\alpha=d_0(\alpha-\hat N\alpha)
    \end{align*}
    and so we get
    \begin{align*}
    d_0 d^{-1}_0|_{\image(\tilde d_0)} &: \image \tilde d_0 \to \image d_0\ ,\ 
    \tilde d_0 \alpha  \longmapsto d_0 A \alpha \ .
    \end{align*}
    The map given by the composition 
    $$\calT \stackrel{A^{-1}}{\longrightarrow} \calT \stackrel{\tilde d_0}{\longrightarrow} \tilde d_0 \calT \stackrel{d_0 d_0^{-1}}{\longrightarrow} \image d_0\ , $$
    simplifies into $d_0\colon\mathcal{T}\to \mathrm{image\,}d_0$, since for any $\alpha\in K$ we have
    $$d_0d_0^{-1}\tilde d_0A^{-1}\alpha=d_0(A -\id +d^{-1}_0 d_0 )A^{-1}\alpha=d_0AA^{-1}\alpha=d_0\alpha.$$
    Since $d_0$ is an invertible map from $\calT$ to $\image d_0$, this implies that $\tilde d_0|_{\calT}$ must be invertible as well. Repeating this argument for every $K_j$ and $(\calT \cap K_j)$, we have $K_j = \ker \tilde d_0|_{K_j} \oplus ( \calT \cap K_j)$. Finally, by construction we have $\mathcal{T}\subset\mathcal{C}$. Moreover the map
    $$K \to \image \tilde d_0, \qquad \alpha \mapsto \tilde d_0 A^{-1} d_0^{-1}\alpha,$$
    will be a projection onto $\image \tilde d_0$ with kernel $\ker d_0^{-1}=\calC$, implying that $\calC$ is transverse to $\image \tilde d_0$.
    \item By (a) and Proposition~\ref{prop:TransInv}, it is sufficient to prove the result in the case where $\calT \subset \calC$ equals $\tilde \calT \subset \tilde \calC$. 
    Continuing with the notation in (a), we observe that $\tilde d_0^{-1} = A^{-1} d_0^{-1}$. Furthermore, if $\tilde P = \id - \tilde d_0^{-1} d - d \tilde d_0^{-1}$, then
    \begin{align*}
    \tilde \calP & = \{ \tilde P(\alpha) = \alpha  \, : \, \alpha \in K \}
    = \{ \alpha \in K \, : \, \tilde d_0^{-1} \alpha = 0 , \tilde d_0^{-1} d\alpha = 0\} \\
    & = \{ \alpha \in K \, : \, A^{-1} d_0^{-1} \alpha = 0 , A^{-1} d_0^{-1} d\alpha = 0\} =  \calP
    \end{align*}
    since $A^{-1}$ is an invertible map of $\calT$ to itself. The result follows.
\end{enumerate}
\end{proof}

\end{document}
